% Part: normal-modal-logic
% Chapter: tableaux
% Section: completeness

\documentclass[../../../include/open-logic-section]{subfiles}

\begin{document}

\olfileid{nml}{tab}{cpl}

\olsection{\Log{K} సంపూర్ణత}

\begin{explain}
  \tetoken{టాబ్లోల}{!!{tableau}s} పద్ధతి సంపూర్ణమని చూపాలంటే,
  $\Gamma \Proves !A$ అని చూపే సంవృత \tetoken{టాబ్లో}{!!{tableau}}
లేనప్పుడల్లా $\Gamma \Entails/ !A$, అంటే ప్రతినమూనా
ఉంటుందని చూపాలి. కానీ ``సంవృత \tetoken{టాబ్లో}{!!{tableau}}
లేదు'' అంటే దాన్ని నిర్మించడానికి ప్రయత్నించే ప్రతి మార్గమూ
సంవృతం కావడంలో విఫలమవుతుంది. ప్రతి మార్గం అలా
విఫలమైతే ఒక నిర్దిష్టమైన \emph{క్రమబద్ధమైన, సమగ్ర}
మార్గమూ విఫలమవుతుందన్నదే కీలకం. సంవృత
\tetoken{టాబ్లో}{!!{tableau}} ఉంటే ఈ క్రమబద్ధమైన,
సమగ్ర మార్గం దాన్ని సంవృతం చేస్తుంది. ఆ ఒక్క టాబ్లోలోని
వివృత శాఖల్లో ప్రతినమూనాను నిర్వచించడానికి కావలసిన
సమాచారమంతా ఉంటుంది. అటువంటి వివృత శాఖ ఇచ్చే
ప్రతినమూనాలో ఆ శాఖపై ఉపయోగించిన పూర్వసూచికలే లోకాలు;
\tetoken{ప్రతిజ్ఞావాక్య చరరాశి}{!!a{propositional
    variable}}~$p$ అనేది~$\sigma$లో సత్యమవుతుంది,
అప్పుడు మాత్రమే $\sFmla{\True}{p}[\sigma]$ ఆ శాఖపై ఉంటుంది.
\end{explain}

\begin{defn}
  \tetoken{ఒక టాబ్లో}{!!a{tableau}} శాఖలో నియమాన్ని
వర్తింపజేయగల పూర్వసూచిక గల \tetoken{సూత్రం}{!!{formula}}
$\sFmla{S}{!A}[\sigma]$ ఉన్నప్పుడల్లా, కింది వాటినీ
కలిగి ఉంటే ఆ శాఖను \emph{సంపూర్ణమైనది} అంటాం:
  \begin{enumerate}
    \item ప్రతిజ్ఞావాక్య వరుస-చేర్పు నియమాలైతే, సంబంధిత
      నిష్కర్షలైన పూర్వసూచిక గల \tetoken{సూత్రాలన్నీ}{!!{formula}s};
    \item ప్రతిజ్ఞావాక్య శాఖావిభజన నియమాలైతే, సంబంధిత
      నిష్కర్ష \tetoken{సూత్రాల్లో}{!!{formula}s} కనీసం ఒకటి;
    \item కొత్త పూర్వసూచిక కావలసిన మోడల్ నియమాలైతే,
      సాధ్యమైన నిష్కర్షల్లో కనీసం ఒకటి;
    \item ఉపయోగించిన పూర్వసూచిక కావలసిన మోడల్ నియమాలైతే,
      శాఖపై కనిపించే ప్రతి అనువర్తనయోగ్య పూర్వసూచికకు
      సంబంధిత నిష్కర్ష.
  \end{enumerate}
\end{defn}

\begin{explain}
ఉదాహరణకు, సంపూర్ణ శాఖలో $\sFmla{\True}{!B \land !C}[\sigma]$
ఉంటే $\sFmla{\True}{!B}[\sigma]$,
$\sFmla{\True}{!C}[\sigma]$ రెండూ ఉంటాయి.
$\sFmla{\True}{!B \lor !C}[\sigma]$ ఉంటే
$\sFmla{\True}{!B}[\sigma]$, $\sFmla{\True}{!C}[\sigma]$
లలో కనీసం ఒకటి ఉంటుంది.
\iftag{prvBox}{$\sFmla{\False}{\Box !B}[\sigma]$}
{$\sFmla{\True}{\Diamond !B}[\sigma]$} ఉంటే, ఏదో ఒక
కొత్త~$n$ కోసం వరుసగా
\iftag{prvBox}{$\sFmla{\False}{!B}[\sigma.n]$}
{$\sFmla{\True}{!B}[\sigma.n]$} కూడా ఉంటుంది.
అలాగే \iftag{prvBox}{$\sFmla{\True}{\Box !B}[\sigma]$}
{$\sFmla{\False}{\Diamond !B}[\sigma]$} ఉంటే,
శాఖపై $\sigma.n$ ఉపయోగించిన ప్రతి~$n$ కోసం వరుసగా
\iftag{prvBox}{$\sFmla{\True}{!B}[\sigma.n]$}
{$\sFmla{\False}{!B}[\sigma.n]$} కూడా ఉంటుంది.
\sourcecorrection{OLTENMLTABCPL-001}{సంపూర్ణ శాఖ ఉదాహరణల్లో
  మూలం సంయోజన పూర్వసూచికను వదిలింది; విడియోజనలో
  సత్యమైన~$!B$ బదులు అసత్యమైన~$!B$ను ఇచ్చింది;
  మోడల్ చిహ్నిత సూత్రాల్లో లోపలి సూత్రాన్ని వదిలి,
  నిష్కర్షల్లో కూడా అసలు సంయోజకాన్నే ఉంచింది.
  ముందటి నియమ పట్టిక ప్రకారం పూర్వసూచికలు,
  సత్య సంకేతాలు,~$!B$ నిష్కర్షలను పునరుద్ధరించాం.}
\end{explain}

\begin{prop}\ollabel{prop:complete-tableau}
  ప్రతి పరిమిత $\Gamma$కు, ప్రతి శాఖా సంపూర్ణంగా ఉండే
  \tetoken{టాబ్లో}{!!a{tableau}} ఉంటుంది.
\end{prop}

\begin{proof}
  $\Gamma$ కోసం ఉన్న \tetoken{ఒక టాబ్లో}{!!a{tableau}}లోని
  వివృత శాఖను పరిశీలిద్దాం. ఆ శాఖలో నియమం వర్తింపజేయగల
  పూర్వసూచిక గల \tetoken{సూత్రాలు}{!!{formula}s}
  పరిమిత సంఖ్యలో ఉన్నాయి. వాటిలో (1)--(4) షరతులు
  ఇంకా నెరవేరని ప్రతి పూర్వసూచిక గల \tetoken{సూత్రానికి}{!!{formula}s}
  ఒక స్థిర క్రమంలో (ఉదాహరణకు, పై నుంచి కిందికి)
  వర్తించే నియమాలను వర్తింపజేసి శాఖను విస్తరించండి.
  కొన్ని సందర్భాల్లో శాఖలు విడిపోతాయి; మిగిలిన
  పూర్వసూచిక గల \tetoken{సూత్రాలకు}{!!{formula}s} ప్రతి
  ఫలిత శాఖ చివర నియమాన్ని వర్తింపజేయండి. శాఖలో
  పూర్వసూచిక గల \tetoken{సూత్రాలు}{!!{formula}s}, ఉపయోగించిన
  పూర్వసూచికలు రెండూ పరిమిత సంఖ్యలోనే ఉన్నందున, ఈ
  ప్రక్రియ చివరకు అసలు శాఖను విస్తరించే (బహుశా అనేక)
  శాఖలను ఇస్తుంది. వాటిలో ప్రతిదానికి ప్రక్రియను
  వర్తింపజేసి పునరావృతం చేయండి. నిర్మాణం ప్రకారం
  చివరకు ప్రతి శాఖ సంపూర్ణంగా ఉంటుంది.
  \sourcecorrection{OLTENMLTABCPL-002}{మూల నిరూపణ చివర
    ప్రతి శాఖ ``సంవృతం'' అవుతుందని ఉంది; ఇది ప్రతిపాదనలోని
    ``సంపూర్ణం'' వాదనకు విరుద్ధం, వివృత శాఖల నుంచే
    ప్రతినమూనా నిర్మించాలనే తదుపరి దశకూ విరుద్ధం.
    ``సంపూర్ణం''గా సరిచేశాం. ఈ పునరావృత ప్రక్రియ
    పరిమిత సమయంలో ముగుస్తుందని మూలం విడిగా నిరూపించదు.}
\end{proof}

\begin{thm}[సంపూర్ణత]
  \ollabel{thm:tableau-completeness}
  $\Gamma$కు సంవృత \tetoken{టాబ్లో}{!!{tableau}} లేకపోతే,
  $\Gamma$ సంతృప్తిపరచదగినది.
\end{thm}

\begin{proof}
పై ప్రతిపాదన ప్రకారం, పరిమితమైన~$\Gamma$కు ప్రతి శాఖా
సంపూర్ణంగా ఉండే \tetoken{టాబ్లో}{!!a{tableau}} ఉంది.
దానికి సంవృత \tetoken{టాబ్లో}{!!{tableau}} లేదు కనుక,
కనీసం ఒక శాఖ వివృతంగానూ సంపూర్ణంగానూ ఉండే
\tetoken{టాబ్లో}{!!a{tableau}} ఉంది. ఆ శాఖపైని పూర్వసూచిక
గల \tetoken{సూత్రాల}{!!{formula}s} సమితి~$\Delta$,
దానిలోని పూర్వసూచికల సమితి~$P(\Delta)$ అనుకుందాం.
\sourcecorrection{OLTENMLTABCPL-003}{పైన ప్రతిపాదన
  పరిమిత~$\Gamma$కు మాత్రమే చెప్పబడింది; మూల సంపూర్ణత
  సిద్ధాంతం, దిగువ పర్యవసాన ఉపసిద్ధాంతాలు పరిమితత్వ
  షరతు లేకుండా ఉన్నాయి. సాధారణ~$\Gamma$కు ఈ అడుగు
  వర్తించాలంటే సంహతత్వం లేదా అనంత శాఖపై న్యాయమైన
  విస్తరణ వంటి అదనపు వాదన కావాలి; మూలంలో అది లేదు.
  ఇక్కడ పరిమిత సందర్భం వరకే నిరూపణ అందుతుందని ప్రకటించాం.}

నమూనా~$\mModel{M(\Delta)} = \tuple{P(\Delta), R, V}$ను
నిర్వచిద్దాం. దీని లోకాలు~$\Delta$లో కనిపించే పూర్వసూచికలు;
ప్రాప్యత సంబంధం ఇది:
\[
R\sigma\sigma' \quad \text{అప్పుడే, అప్పుడు మాత్రమే} \quad
\sigma'=\sigma.n \quad \text{ఏదో~$n$ కోసం}
\]
మరియు
\[
V(p) = \Setabs{\sigma}{\sFmla{\True}{p}[\sigma] \in \Delta}.
\]
ప్రతి పూర్వసూచికను అదే పూర్వసూచిక అనే లోకానికి
కేటాయించే తాదాత్మ్య అర్థనిర్దేశం~$f$ను తీసుకుందాం.
\sourcecorrection{OLTENMLTABCPL-007}{మూలం లోకాలుగా
  పూర్వసూచికలను తీసుకున్నా, ముందటి సంతృప్తి నిర్వచనం
  కోరే అర్థనిర్దేశం~$f$ను స్పష్టంగా చెప్పలేదు;
  పూర్వసూచికలపై తాదాత్మ్య అర్థనిర్దేశాన్ని ప్రకటించి,
  చివరి~$\Gamma$ సంతృప్తి ప్రకటనకు~$[f]$ను చేర్చాం.}
$!A$పై ఆగమనంతో, $\sFmla{\True}{!A}[\sigma] \in \Delta$
అయితే $\mSat{M(\Delta)}{!A}[\sigma]$ అనీ,
$\sFmla{\False}{!A}[\sigma] \in \Delta$ అయితే
$\mSat/{M(\Delta)}{!A}[\sigma]$ అనీ చూపుదాం.

\begin{enumerate}
  \item \indcase{!A}{p}{$\sFmla{\True}{\indfrm}[\sigma] \in \Delta$
    అయితే,~$V$ నిర్వచనం ప్రకారం $\sigma \in V(p)$;
    అందువల్ల $\mSat{M(\Delta)}{\indfrm}[\sigma]$.

    $\sFmla{\False}{\indfrm}[\sigma] \in \Delta$ అయితే
    $\sFmla{\True}{\indfrm}[\sigma] \notin \Delta$; లేకపోతే
    శాఖ సంవృతమవుతుంది. కాబట్టి $\sigma \notin V(p)$;
    అందువల్ల $\mSat/{M(\Delta)}{\indfrm}[\sigma]$.}
  \item \indcase{!A}{\lnot !B}{\iftag{probNot}{వ్యాయామం.}{
      $\sFmla{\True}{\indfrm}[\sigma] \in \Delta$ అయితే,
      శాఖ సంపూర్ణం కాబట్టి
      $\sFmla{\False}{!B}[\sigma] \in \Delta$. ఆగమన
      పరికల్పన ప్రకారం $\mSat/{M(\Delta)}{!B}[\sigma]$;
      అందువల్ల $\mSat{M(\Delta)}{\indfrm}[\sigma]$.

      $\sFmla{\False}{\indfrm}[\sigma] \in \Delta$ అయితే,
      శాఖ సంపూర్ణం కాబట్టి
      $\sFmla{\True}{!B}[\sigma] \in \Delta$. ఆగమన
      పరికల్పన ప్రకారం $\mSat{M(\Delta)}{!B}[\sigma]$;
      అందువల్ల $\mSat/{M(\Delta)}{\indfrm}[\sigma]$.}}
  \item \indcase{!A}{!B \land !C}{\iftag{probAnd}{వ్యాయామం.}{
      $\sFmla{\True}{\indfrm}[\sigma] \in \Delta$ అయితే,
      శాఖ సంపూర్ణం కాబట్టి
      $\sFmla{\True}{!B}[\sigma] \in \Delta$,
      $\sFmla{\True}{!C}[\sigma] \in \Delta$ రెండూ ఉంటాయి.
      ఆగమన పరికల్పన ప్రకారం
      $\mSat{M(\Delta)}{!B}[\sigma]$,
      $\mSat{M(\Delta)}{!C}[\sigma]$; కనుక
      $\mSat{M(\Delta)}{\indfrm}[\sigma]$.

      $\sFmla{\False}{\indfrm}[\sigma] \in \Delta$ అయితే,
      శాఖ సంపూర్ణం కాబట్టి
      $\sFmla{\False}{!B}[\sigma] \in \Delta$ లేదా
      $\sFmla{\False}{!C}[\sigma] \in \Delta$.
      ఆగమన పరికల్పన ప్రకారం
      $\mSat/{M(\Delta)}{!B}[\sigma]$ లేదా
      $\mSat/{M(\Delta)}{!C}[\sigma]$; కనుక
      $\mSat/{M(\Delta)}{\indfrm}[\sigma]$.
      \sourcecorrection{OLTENMLTABCPL-004}{సంయోజనం అసత్యమైన
        శాఖలో రెండో ప్రత్యామ్నాయానికి మూలం పొరపాటుగా
        మళ్లీ~$!B$ అసత్యం అని ఇచ్చింది; శాఖ నియమంలోని
        రెండో ప్రత్యామ్నాయం~$!C$ అసత్యం.}}
  \item \indcase{!A}{!B \lor !C}{\iftag{probOr}{వ్యాయామం.}{
      $\sFmla{\True}{\indfrm}[\sigma] \in \Delta$ అయితే,
      శాఖ సంపూర్ణం కాబట్టి
      $\sFmla{\True}{!B}[\sigma] \in \Delta$ లేదా
      $\sFmla{\True}{!C}[\sigma] \in \Delta$.
      ఆగమన పరికల్పన ప్రకారం
      $\mSat{M(\Delta)}{!B}[\sigma]$ లేదా
      $\mSat{M(\Delta)}{!C}[\sigma]$; కనుక
      $\mSat{M(\Delta)}{\indfrm}[\sigma]$.

      $\sFmla{\False}{\indfrm}[\sigma] \in \Delta$ అయితే,
      శాఖ సంపూర్ణం కాబట్టి
      $\sFmla{\False}{!B}[\sigma] \in \Delta$,
      $\sFmla{\False}{!C}[\sigma] \in \Delta$ రెండూ.
      ఆగమన పరికల్పన ప్రకారం
      $\mSat/{M(\Delta)}{!B}[\sigma]$,
      $\mSat/{M(\Delta)}{!C}[\sigma]$ రెండూ; కనుక
      $\mSat/{M(\Delta)}{\indfrm}[\sigma]$.
      \sourcecorrection{OLTENMLTABCPL-005}{విడియోజనం
        అసత్యమైన శాఖలో రెండో ఆగమన నిష్కర్షకు మూలం
        మళ్లీ~$!B$ అసత్యం అని ఇచ్చింది; రెండో సంయోజ్యం
        $!C$ అసత్యం కావాలి.}}
  \item \indcase{!A}{!B \lif !C}{\iftag{probIf}{వ్యాయామం.}{
      $\sFmla{\True}{\indfrm}[\sigma] \in \Delta$ అయితే,
      శాఖ సంపూర్ణం కాబట్టి
      $\sFmla{\False}{!B}[\sigma] \in \Delta$ లేదా
      $\sFmla{\True}{!C}[\sigma] \in \Delta$.
      ఆగమన పరికల్పన ప్రకారం
      $\mSat/{M(\Delta)}{!B}[\sigma]$ లేదా
      $\mSat{M(\Delta)}{!C}[\sigma]$; కనుక
      $\mSat{M(\Delta)}{\indfrm}[\sigma]$.

      $\sFmla{\False}{\indfrm}[\sigma] \in \Delta$ అయితే,
      శాఖ సంపూర్ణం కాబట్టి
      $\sFmla{\True}{!B}[\sigma] \in \Delta$,
      $\sFmla{\False}{!C}[\sigma] \in \Delta$ రెండూ.
      ఆగమన పరికల్పన ప్రకారం
      $\mSat{M(\Delta)}{!B}[\sigma]$,
      $\mSat/{M(\Delta)}{!C}[\sigma]$ రెండూ; కనుక
      $\mSat/{M(\Delta)}{\indfrm}[\sigma]$.
      \sourcecorrection{OLTENMLTABCPL-006}{షరతు సూత్రం
        అసత్యమైన శాఖలో రెండో ఆగమన నిష్కర్షను మూలం
        $!B$ అసత్యంగా వ్రాసింది; నిష్కర్ష భాగం
        $!C$ అసత్యం కావాలి.}}
  \item \indcase{!A}{\Box !B}{\iftag{probBox}{వ్యాయామం.}{
      $\sFmla{\True}{\indfrm}[\sigma] \in \Delta$ అయితే,
      శాఖ సంపూర్ణం కాబట్టి, శాఖపై ఉపయోగించిన ప్రతి
      $\sigma.n$కు $\sFmla{\True}{!B}[\sigma.n] \in \Delta$;
      అంటే $R\sigma\sigma'$ అయ్యే ప్రతి
      $\sigma' \in P(\Delta)$కు ఇదే వర్తిస్తుంది. ఆగమన
      పరికల్పన ప్రకారం $R\sigma\sigma'$ అయ్యే ప్రతి~$\sigma'$కు
      $\mSat{M(\Delta)}{!B}[\sigma']$. కనుక
      $\mSat{M(\Delta)}{\indfrm}[\sigma]$.

      $\sFmla{\False}{\indfrm}[\sigma] \in \Delta$ అయితే,
      శాఖ సంపూర్ణం కాబట్టి ఏదో~$\sigma.n$కు
      $\sFmla{\False}{!B}[\sigma.n] \in \Delta$.
      ఆగమన పరికల్పన ప్రకారం
      $\mSat/{M(\Delta)}{!B}[\sigma.n]$.
      $R\sigma(\sigma.n)$ కాబట్టి, అటువంటి
      $\mSat/{M(\Delta)}{!B}[\sigma']$ అయ్యే
      $\sigma'$ ఉంది. కనుక
      $\mSat/{M(\Delta)}{\indfrm}[\sigma]$.}}
  \item \indcase{!A}{\Diamond !B}{\iftag{probDiamond}{వ్యాయామం.}{
      $\sFmla{\True}{\indfrm}[\sigma] \in \Delta$ అయితే,
      శాఖ సంపూర్ణం కాబట్టి ఏదో~$\sigma.n$కు
      $\sFmla{\True}{!B}[\sigma.n] \in \Delta$.
      ఆగమన పరికల్పన ప్రకారం
      $\mSat{M(\Delta)}{!B}[\sigma.n]$.
      $R\sigma(\sigma.n)$ కాబట్టి, అటువంటి
      $\mSat{M(\Delta)}{!B}[\sigma']$ అయ్యే
      $\sigma'$ ఉంది. కనుక
      $\mSat{M(\Delta)}{\indfrm}[\sigma]$.

      $\sFmla{\False}{\indfrm}[\sigma] \in \Delta$ అయితే,
      శాఖ సంపూర్ణం కాబట్టి, శాఖపై ఉపయోగించిన ప్రతి
      $\sigma.n$కు $\sFmla{\False}{!B}[\sigma.n] \in \Delta$;
      అంటే $R\sigma\sigma'$ అయ్యే ప్రతి
      $\sigma' \in P(\Delta)$కు ఇదే వర్తిస్తుంది. ఆగమన
      పరికల్పన ప్రకారం $R\sigma\sigma'$ అయ్యే ప్రతి~$\sigma'$కు
      $\mSat/{M(\Delta)}{!B}[\sigma']$. కనుక
      $\mSat/{M(\Delta)}{\indfrm}[\sigma]$.}}
\end{enumerate}
$\Gamma \subseteq \Delta$ కాబట్టి,
$\mSat{M(\Delta)}{\Gamma}[f]$.
\end{proof}

\begin{prob}
\olref[nml][tab][cpl]{thm:tableau-completeness} నిరూపణను పూర్తి చేయండి.
\end{prob}

\begin{cor}
\ollabel{cor:entailment-completeness}
$\Gamma \Entails !A$ అయితే $\Gamma \Proves !A$.
\end{cor}

\begin{cor}
\ollabel{cor:weak-completeness}
$!A$ అన్ని నమూనాల్లో సత్యమైతే, $\Proves !A$.
\end{cor}

\end{document}
